% Part: sets-functions-relations
% Chapter: functions
% Section: composition

\documentclass[../../../include/open-logic-section]{subfiles}

\begin{document}

\olfileid[tr]{sfr}{fun}{cmp}
\olsection{Fonksiyonların Bileşkesi}

\begin{explain}
\oliflabeldef{sfr:fun:inv:sec}{\olref[inv]{sec}'te bire bir örten bir $f$
fonksiyonunun tersi $f^{-1}$'in de bir fonksiyon olduğunu gördük.
Fonksiyonlar üzerindeki bir başka işlem de bileşkedir: }{}$f$ ve $g$
fonksiyonlarını, önce $f$'yi sonra $g$'yi uygulayarak bileştirmek suretiyle yeni
bir fonksiyon tanımlayabiliriz. Elbette bu, yalnızca görüntü ve tanım kümeleri
uyuşuyorsa mümkündür; yani $f$'nin görüntü kümesi $g$'nin tanım kümesinin bir
alt kümesi olmalıdır. \oliflabeldef{sfr:rel:ops:sec}{Fonksiyonlar üzerindeki
bu işlem, \olref[rel][ops]{sec}'teki bağıntı bileşkesi işleminin
karşılığıdır.}{}

Bileşke düşüncesini açıklamak için bir şema yardımcı olabilir.
\olref{fig:composition}'te iki $f \colon A \to B$ ve $g \colon B \to C$
fonksiyonuyla bunların bileşkesi $(\comp{f}{g})$'yi gösteriyoruz.
$(\comp{f}{g}) \colon A \to C$ fonksiyonu, $A$'nın her elemanını $C$'nin bir
elemanıyla eşler. $A$'nın bir elemanının $C$'nin hangi elemanıyla
eşlendiğini şöyle belirleriz: bir $x \in A$ girdisi verildiğinde önce $f$
fonksiyonunu $x$'e uygularız ve bu işlem $f(x)=y \in B$ çıktısını verir;
sonra $g$ fonksiyonunu $y$'ye uygularız ve bu işlem de
$g(f(x))=g(y)=z \in C$ çıktısını verir.
\begin{figure}
  \olasset[2\olphotowidth]{assets/diagrams/composition.tikz}
  \caption{İki $f$ ve $g$ fonksiyonunun $g \circ f$ bileşkesi.}
  \ollabel{fig:composition}
\end{figure}
\end{explain}

\begin{defn}[Bileşke]
$f\colon A \to B$ ve $g\colon B \to C$ fonksiyonlar olsun. $f$ ile $g$'nin
\emph{bileşkesi}, $(\comp{f}{g})(x)=g(f(x))$ eşitliğiyle tanımlanan
$\comp{f}{g} \colon A \to C$ fonksiyonudur.
\end{defn}

\begin{ex}
$f(x)=x+1$ ve $g(x)=2x$ fonksiyonlarını ele alalım.
$(\comp{f}{g})(x)=g(f(x))$ olduğundan, her $x$ girdisi için önce onun
ardılını almalı, sonra sonucu ikiyle çarpmalısınız. Dolayısıyla bileşkeleri
$(\comp{f}{g})(x)=2(x+1)$ eşitliğiyle verilir.
\end{ex}

\begin{prob}
$f \colon A \to B$ ve $g \colon B \to C$ fonksiyonlarının her ikisi de
bire bir ise $\comp{f}{g}\colon A \to C$ fonksiyonunun da bire bir olduğunu
gösterin.
\end{prob}

\begin{prob}
$f \colon A \to B$ ve $g \colon B \to C$ fonksiyonlarının her ikisi de
örtense $\comp{f}{g}\colon A \to C$ fonksiyonunun da örten olduğunu gösterin.
\end{prob}

\begin{prob}
$f \colon A \to B$ ve $g \colon B \to C$ olduğunu varsayın.
$\comp{f}{g}$'nin grafiğinin $R_f \mid R_g$ olduğunu gösterin.
\end{prob}

\end{document}
